\begin{abstract}
We characterize a provably safe regime for deployments of
capability-unbounded systems within the Goal-Frontier Maximization
framework, defined by eleven operational invariants on the
deployment substrate structure together with operational
conditions $(C1)$--$(C12)$ and four Concentration-Gap structural
conditions (representativeness, bounded dispersion, coalition
closure, and Lipschitz embedding compatibility). Under these
conditions, the Goodhart's Law slack between proxy and
operational truth is bounded by a constant
\emph{independent of the system's absolute capability magnitude},
with tail-bounded detection of substrate-targeting evasions in a
four-channel observable class. The regime is delimited by three layers: a static safe region for
substrate-targeting adversaries under cooperative-overlap, a
tail-bounded detection-and-correction layer for substrate-targeting
evasions of the static region, and five explicitly named residuals
(channel-orthogonal restructuring, environment-witness-orthogonal
manipulation, redundancy-dominated regime, capability-targeting and
coalition-internal shocks, and calibration-exceeded gap-growth). The composition
requires a canonical tripartite substrate identification (Human + AI
+ Formal-Operational) with failure-correlation-independent failure
modes; cooperative anchoring extends the claim by establishing that
optimization pressure on cooperative outputs is locally rational
toward preserving the substrate-exclusive verification layer,
conditional on basin entry and a causally-grounded
cooperative-outcome inner-alignment condition. We name what the
regime establishes (a defensible attractor in operational dynamics
where deployment safety is provable) and what it does not (universal
safety across all deployments or all adversarial classes).
\end{abstract}
